%----------------------------------------------------------------------
\chapter{Software components detailed design}
\label{ch:detailed}
%----------------------------------------------------------------------

This chapter describes the detailed implementation of the different
software components.

%----------------------------------------------------------------------
\section{Challenges}
\label{sc:challenges}
%----------------------------------------------------------------------

Static stack analysis techniques have a common set of challenges that
need to be addressed:

\begin{itemize}

\item {\it Cycles in the Control Flow Graph}:

In presence of control flow cycles involving stack consumption, the worst case
bound is a function of the maximum number of cycle iterations.

When bounds on such iteration counts are unknown, the amount of potential
stack usage is infinite.

Bounding the iteration count is very often difficult or impossible.

\item {\it Indirect jumps or subprogram calls}:

Determining where such transfers of control lead is often not achievable,
and the corresponding stack usage is then unknown.

\item
{\it Unresolved calls}:

These are calls to subprograms for which no stack usage information is
available, introducing an unknown amount of stack usage in a call chain.

\item {\it Dynamic stack allocations}:

For instance from {\tt alloca} with a variable argument in C, or from
dynamically sized local objects in Ada.

They introduce potentially difficult to bound variable amounts of stack usage
on the paths where they appear.

\item {\it Possible interrupt events}:

When interrupt (hardware or signal) handlers run on the currently active
stack, their consumption has to be accounted for when sizing the various areas
in which they may run.

When they run on their own dedicated stack, this one has to be properly sized
too.

In any case, how interrupts may preempt each other greatly affects the
associated possible stack usage and is not easy to determine automatically.

\cite{pal:icse-01, pal:ftrtft-02, pal:sas-03} and
\cite{reg:emsoft-03} are examples of research literature on this matter.

\item {\it Dynamic global behavior}:

Analysis tools typically include in their evaluations many paths that can
never be taken at run-time, referenced as False Paths in \cite{bry:thesis} or
\cite{alt:erts-96}.

Excluding some of such paths on safe grounds may, for instance, allow cutting
some control flow cycles or the production of tighter worst case bounds,
saving run-time resources.

\cite {reg:emsoft-03} illustrates the use of such techniques to automatically
compute the hardware interrupt preemption graph for an AVR target.

This is a hard problem in the general case.

\end{itemize}

It will be shown in the following sections where and how these
challenges are addressed.

%----------------------------------------------------------------------
\section{Compiler}
%----------------------------------------------------------------------

We are only going to describe the implementation of the two new
command line options used for static stack analysis. The options are
independent from each other, although we need both in order to get 

\subsubsection{Stack usage}

The general principle is as follows: we directly leverage the long existing
infrastructure for calculating the frame size of functions, made up of both
a generic part and a target back-end dependent part, to report the amount of
static stack usage for the function being compiled.

As the back-end dependent part already needs to gather the final result of
the calculation before emitting the assembly code, the actual implementation
essentially boils down to modifying every back-end so as to be able to easily
retrieve this result by means of a ``target hook''.

We found that, at least for the most common architectures, the changes to
be made are very localized.

The ``target hook'' implementation model proved to be a very efficient device
and really minimizes the effort required to add support for a new target.

x86, powerpc, sparc, alpha, mips and hppa have been covered up to now.

The only challenge is to make sure that every single byte allocated on the
stack by the calling sequence, even if it is not formally part of the frame,
is taken into account.

In this context, one interesting technical point is of note: the difference in
treatment between \verb|ACCUMULATE_OUTGOING_ARGS| and \verb|PUSH_ARGS|
targets.

In the former case, the arguments of called functions are accounted for in the
final frame size, whereas they are not in the latter case; moreover, another
subtlety comes into play in the latter case, in the form of the {\tt
-fdefer-pop} command line option, which instructs the compiler to not pop the
pushed arguments off the stack immediately after the call returns.

This may result in increased stack usage and requires a special circuitry
to be properly dealt with at compile-time.

The remaining task is then to detect the dynamic stack usage patterns, much
like what is implemented to support {\tt -fstack-check}.

For this initial implementation, we mainly operate in the Tree to RTL expander
by intercepting the requests for dynamic stack allocation.

Moreover, as we are primarily interested in bounding them, we also try to
deduce from the IL (here RTL) limits easily evaluated at compile-time.

A technical detail that must not be overlooked here is that the compiler
may generate code to dynamically align the stack pointer.

While the amount of stack usage is easily bounded in that case, it must not be
forgotten in the final result.

There is certainly room for improvement in either direction on the axis of
compiler passes: by working later down the RTL optimization passes, one
should be able to obtain additional constant bounds for dynamic allocation
cases that are not inherently bounded;

by working closer to the trees handed down by the front-end, one should be
able to recognize inherently static allocation patterns that happen to require
a dynamic-like treatment for specific reasons, as is the case for big objects
when \prog{-fstack-check} is enabled for example.

\subsubsection{Call graph}

We again leverage an existing infrastructure of the compiler to implement
this command line option, but a much more recent one, that is the "callgraph"
machinery introduced in the 3.4.x series.

We only use this machinery to gather information all the way through the
compilation.

We don't drive the compilation with it or attempt to rely on it for any code
generation purpose, so our approach is not tied to one of the compilation
modes supported by GCC 3.4.x.

In particular, it works with the unit-at-a-time mode (C and C++ compilers at
-O2), with the function-at-a-time mode (C and C++ compilers at -O0/-O1) and
finally the statement-at-a-time mode (Ada compiler).

The general principle is straightforward: we record every direct function
call the compiler has to process, either at the Tree level in unit-at-a-time
mode or at the RTL level in non unit-at-a-time modes, and every indirect
function call at the RTL level.

Of course some of these function calls may end up being optimized away at one
stage of the compilation but, as we aim at computing a worst case scenario,
this conservative stance is appropriate.

However, an optimization technique relating to function calls is particularly
of note since it can bring about huge differences in results for any types of
callgraph-based analysis, depending on whether it is accounted for or not,
that is function inlining.

We therefore do arrange for eliminating or not registering in the first place
in the call graph edges that correspond to function calls for which the callee
is inlined in the caller in the assembly code emitted by the compiler.

For the example of the {\fsu} option, the immediate benefit is that the
static stack usage of the callee is guaranteed not to be counted twice
in the final calculation.

When additional analysis-oriented command line options are passed to the
compiler, nodes in the call graph can be annotated with additional
information gathered all the way through the compilation.

The final step is to generate an output suited to the kind of post-processing
callgraph-based analyzes are naturally likely to require.

In particular, as the process of merging .ci files is somewhat akin to linking
object files, the problem of the uniqueness of the identifiers for the merge
points has to be addressed.

The VCG format for the output was chosen as a simple and natural medium
to convey compiler-generated, callgraph-based information.

It goes without saying that targeting any other (text-based) format is easily
doable and would only command modifications at the very end of the chain.


%----------------------------------------------------------------------
\section{Call graph generator}
%----------------------------------------------------------------------


%----------------------------------------------------------------------
\subsection{Cycles}
%----------------------------------------------------------------------

One of the problems identified in section \ref{sc:challenges} is the
detection and handling of potential cycles in the control flow. There
are three basic aspects:

\begin{itemize}

\item How to detect a cycle.

\item How to represent, and identify uniquely, a cycle. A cycle is a set of nodes

\item How to characterize stack usage for a cycle. We need an
  attribute representing the maximum number 

\end{itemize}

It is important to detect these cycles.

We could thing about using an external tool for bounding these cycles
using control and data flow analysis. The user can also bound these
cycles manually.

%----------------------------------------------------------------------
\subsection{Indirect calls (dereference of access-to-subprogram)}
%----------------------------------------------------------------------

One challenge defined in section \ref{sc:challenges} has to do
with indirect calls in the form of a dereference of
access-to-subprogram. When the compiler detects such a call, it
generates a \emph{dummy proxy} node into the call graph
information. In order to fill the information about possible targets
for this call, we could analyze all the subprograms in the
partition. All the subprograms to which a \prog{'Access} or
\prog{'Unrestricted\_Access} apply, and with a signature conforming to
the access value whose dereference generates the indirect call, are
potential candidate targets for the indirect call under analysis.

The user can either provide the missing information about the set of
possible target 

%----------------------------------------------------------------------
\subsection{Dispatching calls}
%----------------------------------------------------------------------

Another form of indirect calls is dispatching calls. challenge defined
in section \ref{sc:challenges} has to do with indirect calls into the
call graph information. When the compiler detects such a call, it
generates a \emph{dummy proxy} node. In order to fill the information
about possible targets for this call, we could analyze all the methods
that could be the target of the dispatching call being analyzed.

The user can either provide the missing information about the set of
possible target

%----------------------------------------------------------------------
\subsection{Unresolved calls}
%----------------------------------------------------------------------

TBC

%----------------------------------------------------------------------
\section{Call graph parser}
%----------------------------------------------------------------------

%----------------------------------------------------------------------
\subsection{Dynamic frames}
%----------------------------------------------------------------------

The user will be able to specify
TBD

%----------------------------------------------------------------------
\section{Worst case analyzer}
%----------------------------------------------------------------------

The result of the worst case analyzer will contain markers that will
indictate how fiable are the results.

There are three major cases:

\begin{itemize}

\item Value computed using only information generated by the
  compiler. Reliable.

\item Value computed using information provided by the user. Less reliable.

\item Values computed using default values. Even less reliable.

\end{itemize}

TBD

%----------------------------------------------------------------------
\section{Modes of operation}
%----------------------------------------------------------------------

There are two major modes of operation: safe and relaxed mode

The safe mode is all about automatically computed bounds. There will
be warnings for every undefined entity (unresolved or indirect calls,
user-defined information, unbounded stack frames, etc.)

The relaxed mode is more flexible. User-defined values are allowed,
and the user may also use defaults for maximum recursion depth, 

   --safe-mode
   --relaxed-mode

   --recursion\_depth[=1]
   --dynamic\_frames[=0]
   --unknown\_call[=0]

The tool will automatically generate skeletons for the missing parts
in the call graph (unresolved or indirect calls, user-defined
information, unbounded stack frames, etc.)

   --generate-ciu
   --overwrite-ciu

%----------------------------------------------------------------------
\section{GNAT Programming Studio}
%----------------------------------------------------------------------

TBD
